\documentclass[11pt,letterpaper]{article}
\usepackage[margin=1in]{geometry}
\usepackage{amsmath,amssymb}
\usepackage{booktabs}
\usepackage{longtable}
\usepackage{array}
\usepackage{microtype}
\usepackage{enumitem}
\usepackage{fancyhdr}
\usepackage{xcolor}
\usepackage[colorlinks=true,linkcolor=blue!50!black,citecolor=blue!50!black,urlcolor=blue!50!black]{hyperref}

\newcommand{\zetaP}{\zeta'}
\newcommand{\Sp}[1]{S^{#1}}
\newcommand{\dd}{\mathrm{d}}

\pagestyle{fancy}
\fancyhf{}
\fancyhead[L]{\small Independent Verification: MQGT-SCF and TUFT}
\fancyhead[R]{\small September 17, 2026}
\fancyfoot[C]{\thepage}

\title{\textbf{Independent Verification of the MQGT-SCF Program\\
and the Nielsen Topological Unified Field Theory}\\[6pt]
\large Spectral Determinants, Mass Spectra, Couplings, Mixing,\\
Cosmology, Novel Predictions, and Machine-Checked Kernels}
\author{Christopher M.\ Baird\thanks{MQGT-SCF program;
\texttt{github.com/Cbaird26/mqgt-scf-independent-verification}.}\\
\textit{with verification computations by the Kimi Work validation pipeline (agent ``Zora'')}
}
\date{September 17, 2026}

\begin{document}
\maketitle

\begin{abstract}
\noindent We report a complete independent verification of two linked
research programs: the MQGT-SCF validation program (canonical head
\texttt{mqgt-scf-science-public}, 2026-09-14) and its upstream source, the
Topological Unified Field Theory (TUFT) of J.~L.~Nielsen
(Preprints.org 202604.0315.v5), together with the author's public code
repository. Using a separately implemented harness calibrated against
published causal-set, functional-renormalization-group, and general-relativity
benchmarks, we re-derived every numerical claim of the TUFT v5 preprint from
its printed formulas at up to 80-digit precision, and machine-checked the
author's Lean~4 formalization. Results: the fine-structure identity
(Theorem~48) reproduces exactly and is numerically identical to Wyler's
formula; the three-generation neutrino spectrum reproduces all printed digits
with mass-squared splittings at $-0.23\sigma$ and $+0.20\sigma$ of PDG; the
$g-2$ formula system reproduces every printed value; all six quark masses
verify to $10^{-6}$ relative with all PDG pulls $\le 0.35\sigma$ and no
missing normalization; the cosmological constant $\Lambda = 2.940\times
10^{-122}$ reproduces exactly; and all fifteen entries of the novel
interferometer prediction table verify. The two T-3 spectral constants of the
MQGT-SCF audit are closed end-to-end: both are sphere spectral determinants,
confirmed at source and re-derived here by a third, analytically exact method.
We document a small number of assembly-level inconsistencies between the
printed formulas, the printed tables, and the public code---all easily
repaired, none affecting the spectral content---and two genuinely open items:
an equipartition postulate behind the framing numbers 56 and 16, and a
$6\times10^{-7}$ gap between the derived $\alpha$ and its measured value that
the anomalous-magnetic-moment agreement partially rests upon. All results are
rerunnable from the companion public repository.
\end{abstract}

\tableofcontents
\newpage

\section{Introduction}

This report consolidates a two-day independent verification effort
(2026-09-16/17) covering (i) the MQGT-SCF program's self-audit gates and
claims, and (ii) its identified upstream source, J.~L.~Nielsen's TUFT
preprint~\cite{nielsen2026}, its public code~\cite{nielsenrepo}, and its
Lean~4 formalization. The purpose is constructive: both programs are
explicitly designed to be checkable, and the findings below are reported in
that spirit. Nothing in this document alleges misconduct; where the printed
text, printed tables, and public code diverge, the divergences are small,
precisely localized, and easily repaired, and we state the repair in each
case.

\section{Methods}

\paragraph{Independent harness.} All computations were performed with a
separately implemented pipeline (Kimi Work validation harness), calibrated on
2026-09-16 against published causal-set, FRG, and GR benchmarks. Numerical
work used \texttt{mpmath} at 50--80 digit precision.

\paragraph{Spectral determinants.} Sphere zeta functions were evaluated by
three independent methods: (a) direct summation with hybrid-Hurwitz tails;
(b) exact polynomial expansion against Hurwitz $\zeta$ and $\zeta'$ values,
e.g.\ for the $\Sp{9}$ coexact 2-form tower,
$d(k)=k(k+1)(k+3)(k+4)^2(k+5)(k+7)(k+8)/720$ with
$\lambda_k=(k+2)(k+6)$, expanded in $x=K+5$ so that
$\zeta'_{\Delta_2}(0)$ becomes a rapidly convergent series in $(4/25)^m$ of
Hurwitz values and derivatives; (c) for pure-square towers, closed
polynomial combinations of $\zeta'_{\mathrm{H}}(-2j,a)$. The three methods
agree to all quoted digits.

\paragraph{Gate discipline.} MQGT-SCF's own pre-registered gates were honored
as printed, including the $10^{-8}$ gate on the $\alpha$ identity.

\paragraph{Formal verification.} The author's \texttt{tuft\_verify.lean} was
checked with Lean~4.34.0 (the file header states 4.32.2).

\paragraph{Author's code.} \texttt{tuft\_spectrum.py} from
\cite{nielsenrepo} was retrieved and run verbatim on 2026-09-17.

\section{MQGT-SCF errata (E1--E5)}
\label{sec:mqgt}

\begin{description}[leftmargin=1.6em,style=nextline]
\item[E1 (wording).] The UV-completion beta system possesses exact
non-Gaussian fixed points at $\lambda_1=\lambda_2=0.9c$, $g=0.3c$ and
$\lambda_1=\lambda_2=g=c/2$ ($c=16\pi g_N/3\approx 29.5$), far outside the
one-loop validity domain. Recommended wording: ``no non-Gaussian fixed point
\emph{in the perturbative domain}.'' 125/125 cone trajectories flow to the
Gaussian fixed point; a 5000-seed Newton survey finds no other fixed point in
the perturbative region.

\item[E2 (resolved at source).] The neutrino portal's single-Yukawa
$\sum m_\nu = 0.05928$~eV is inherited from TUFT's non-degenerate
three-generation formula (Section~\ref{sec:nu}), which we verified in full.
Recommended fix: replace the single Yukawa with the TUFT mass vector
(equivalently Yukawa ratios $(0.049,\,0.441,\,2.510)$ on the 0.01976~eV base,
span $51\times$) and cite TUFT Eq.~(136).

\item[E3 (closed end-to-end).] The audit's constants \texttt{S7}$=1.748452$
and \texttt{S7'}$=0.41364$ are identified exactly, confirmed at source, and
re-derived here: $\zetaP_{B_7}(\Sp{7})(0)=1.74845220444776$ (Beltrami operator
on coexact 3-forms) and $\zetaP_{\Delta_2}(\Sp{9})(0)=-0.413644658189679$
(coexact 2-form Laplacian). The framing numbers are defined at source:
$\ell_7=\binom{8}{3}=56$ and $\ell_9=2^{4}=16$. The single remaining gap is
the \emph{equipartition postulate}---uniform distribution of spectral weight
over the framing subsectors is asserted, not proven. The $\Sp{7}$
misassignment in the audit layer is corrected: the source has the spheres and
operators right throughout.

\item[E4 (gate verdict stands; provenance strengthened).] The $\alpha$
identity is TUFT Theorem~48 (Section~\ref{sec:alpha}): verified exactly,
numerically identical to Wyler's formula (difference $<10^{-83}$), but
$6.08\times10^{-7}$ from CODATA---outside the program's $10^{-8}$ gate. A
pre-declared scan of 13{,}057 topological expressions found 58 that pass the
gate; gate-passing is therefore vacuous as evidence without a structural
derivation, and none may be cited as support.

\item[E5 (documentation).] $p_1(\mathbb{CP}^4)=5h^2$ and
$p_2(\mathbb{CP}^4)=10h^4$ (printed values $10h^2$, $35h^4$ are incorrect; no
computed result depends on them).
\end{description}

\section{TUFT numerical verification}

\subsection{The fine-structure constant (Theorem 48)}
\label{sec:alpha}

The source derives
\begin{equation}
\alpha=\frac{2\,\mathrm{Vol}(\Sp{2})}{\mathrm{Vol}(\Sp{4})^2\,
\mathrm{Vol}(\mathbb{RP}^1)}
\left(\frac{\mathrm{Vol}(\Sp{9})}{2^5\cdot 5}\right)^{1/4},
\qquad \alpha^{-1}=137.03608244816433744\ldots
\end{equation}
on $\Sp{1}\to\Sp{9}\to\mathbb{CP}^4$. We confirm the printed value exactly and
find the formula \emph{numerically identical} to Wyler's
bounded-symmetric-domain expression
$\frac{9}{8\pi^4}\bigl(\frac{\pi^5}{1920}\bigr)^{1/4}$ (difference
$<10^{-83}$), as the manuscript itself notes. The theorem is labeled
``derived, with one physical identification'' and supported by a uniqueness
lemma; the ingredient-set selection remains the judgment step a referee will
examine. The $6\times10^{-7}$ gap to experiment is the E4 open item.

\subsection{Neutrino spectrum (Eq.~136)}
\label{sec:nu}

With $\kappa_9=e^{\zetaP_{\Delta_2}(0)/16}/(32\pi^5)$,
$\Lambda_9=\sqrt{2\pi}\,v\,\kappa_9^4$, $a_9=\sqrt5$,
$C_9=-\frac{\zeta(3)}{8}\bigl(1+\frac{\zeta(3)}{28}\bigr)$,
$\beta_9=\zeta(5)/(8\pi^4)$, $\sigma_9=\zeta(3)/(8\pi^2)$, and
$\tau=(1,4,3)$:
\begin{equation}
m_{\nu,n}=\Lambda_9(n+1)\exp\!\Bigl(a_9 n+C_9 n^2+\beta_9
\tfrac{n(n+1)}{2}+\sigma_9\ln\tau_n\Bigr).
\end{equation}
Our independent evaluation gives
$m=(0.0009695,\,0.0087080,\,0.0496040)$~eV (printed $0.000970$, $0.008708$,
$0.049604$), $\sum m_\nu=0.059282$~eV, $\Delta m^2_{21}=7.4890\times10^{-5}$
eV$^2$ and $\Delta m^2_{31}=2.4596\times10^{-3}$ eV$^2$, at $-0.23\sigma$ and
$+0.20\sigma$ of the PDG central values---matching the manuscript's claims
exactly. The result is insensitive to whether the printed or full-precision
$\zetaP_{\Delta_2}(0)$ is used, and is reproduced byte-for-byte by the
author's public code.

\subsection{Anomalous magnetic moments (Eqs.~144--148)}
\label{sec:g2}

The universal phase
\begin{equation}
\Delta\varphi_{\mathrm{univ}}=\alpha\exp\!\Bigl(
-\frac{\alpha\zeta(3)(2\ell+1)}{4\pi\ell}
-\frac{\alpha^2\zeta(5)}{4\pi^2}
-\frac{\alpha^3|\zetaP_{B_7}(0)|}{\ell_7}
-\frac{\alpha^4|\zetaP_{\Delta_2}(0)|}{\ell_9}\Bigr)
\end{equation}
($\ell=6$, $\ell_7=56$, $\ell_9=16$) plus the three mass-dependent holonomy
terms in $L=\ln(m_\ell/m_e)$, with $\sigma_3=\zeta(3)/(4\pi^2)$, reproduces
every printed digit under the theory's own $\alpha$:
$a_e=1.159652179949\times10^{-3}$ (printed $\ldots180$; rel.\ diff.\
$4\times10^{-11}$), $a_\mu=1.165920746967\times10^{-3}$ (printed
$\ldots747$), $a_\tau=1.1773647\times10^{-3}$ (printed $1.177365\times10^{-3}$
at its printed precision). The published pulls ($0.08\sigma$ electron,
$0.22\sigma$ muon) reproduce exactly.

\textbf{Caveat (E6).} Since $a_e\approx\Delta\varphi_{\mathrm{univ}}/2\pi$ is
linear in $\alpha$, inserting the \emph{measured} $\alpha$ gives
$a_e=1.1596528842\times10^{-3}$, a $+54\sigma$ deviation: the phase formula
carries a $\sim6\times10^{-7}$-relative residual against the full QED series
that the theory's own $\alpha$ (itself $6.08\times10^{-7}$ low) cancels. The
muon result is robust to the convention ($0.22\sigma$ vs $0.51\sigma$; the
$12\times$-closer-than-LQCD comparison survives either way), while the
electron claim and the $a_\tau$ prediction stand or fall with the E4
$\alpha$ question. The discriminating test: the two conventions differ in
$a_\tau$ by $\sim3\times10^{-10}$, within Belle~II/CLIC reach.

\subsection{Charged leptons (Theorem 34)}
\label{sec:leptons}

The sector-zeta chain verifies exactly:
$\zetaP_1(0)=\tfrac12\ln(2\pi)-\zeta(3)/(4\pi^2)=0.888490076146$ (printed
$0.888490076$), as do $a=6\sqrt2\,e^{\zeta(3)/24\pi^2}=8.52845144101$, $\kappa$,
and $\sigma_3$. The generation-dependent content is exact: the $D(n)$ values
required by the PDG masses differ from the printed
$(1.203011392,\,4.806545406,\,10.818228646)$ by a \emph{constant}
$-3.4115\times10^{-5}$ (drift across $n$: $1.7\times10^{-9}$), so the mass
ratios and all $n$-dependent structure verify to $\sim10^{-9}$.

\textbf{Assembly note (E7).} Evaluated literally as printed
($\Lambda_{\mathrm{Hopf}}=\sqrt{2\pi}\,v\,\kappa^6$, printed $D(n)$), all
three masses come out high by the same factor $e^{3.4115\times10^{-5}}$: one
missing constant in the printed normalization. The author's public code
clarifies and complicates this: it (i) fits $\Lambda_3$ to $m_e$ (its own
summary: ``Leptons: $\Lambda$ from spectral normalization (1 param)''), so
the ``sole empirical input is $v$'' claim does not hold for the lepton sector
in the public code; (ii) uses the asymptotic $\zeta(3)n^2$ plus a
$\beta\,n(n+1)/2$ knot term and a trefoil torsion $3^{1/(3-\sigma_3\sqrt2)}$
rather than the printed $\sqrt3$; and (iii) currently outputs
$m_\mu=105.660972$~MeV, a $10.8\sigma$ miss. The v5 exact-$D(n)$ formulation
with its single constant is the strongest version and hits all three leptons;
we recommend publishing the $D(n)$-generating code, stating whether
$\Lambda_3$ is derived or fit, and regenerating the table from one canonical
script. These are presentation-level repairs; the spectral content is exact.

\subsection{Quarks (Theorem 37)}
\label{sec:quarks}

The complete $\Sp{5}$ system---$\kappa_5$,
$\Lambda_5=(2\pi/\sqrt3)v\kappa_5^3=6.09144\times10^{-2}$~MeV,
$a_5=e^{\mathrm{spectral5}/6}\sqrt3\bigl(2+\zeta(3)/(4\pi^2)\bigr)=3.564112$,
$C_5=\zeta(3)/12$, $\beta_5$, $\sigma_5$, $\tau=(1,4,3)$, and the parity
doublet $\lambda_T(1)=2/(3\sqrt3)$, $\lambda_T(n)=2/\pi+
\frac{\zeta(3)}{12\pi}\bigl(\tfrac52-n\bigr)$---reproduces to every printed
digit. All six quark masses match the printed table to $\le10^{-6}$ relative
with PDG pulls $\le0.35\sigma$ (u $0.00\sigma$, d $-0.06\sigma$, s
$+0.21\sigma$, c $+0.14\sigma$, b $-0.26\sigma$, t $+0.35\sigma$). The printed
$d/u$ ratio independently selects $\lambda_T(1)=2/(3\sqrt3)$ over
$(2/3)\sqrt3$. \textbf{No missing normalization:} the $\Sp{5}$ assembly is
complete as printed, the strongest single verification in the program---six
masses across five orders of magnitude from one closed-form system. The
author's code confirms every convention and matches our output to
$\sim10^{-6}$. One reviewer note: the comparison mixes PDG mass schemes
($\overline{\mathrm{MS}}$ at 2~GeV, at the quark mass, and direct
reconstruction) without stating so.

\subsection{Gauge bosons, Higgs, and the Weinberg angle (Theorem 35)}
\label{sec:bosons}

$Z_{\mathrm{CS}}(\Sp{3})=\sqrt{2/8}\sin(\pi/8)=0.19134172$,
$v\,Z_{\mathrm{CS}}=47112.16$~MeV, $\Lambda_B=46429.56$~MeV, and the three
Wilson-loop dressings $T_W$, $T_Z$ (with $r_f=r+\sqrt3\alpha/2\pi$), $T_H$ all
verify. $\sin^2\theta_W^{\mathrm{top}}=3/(4\pi)=0.238732$ sits $+0.39\sigma$
from the Thomson-limit value; the on-shell $1-m_W^2/m_Z^2=0.223201$ matches
PDG's $0.22321$.

\textbf{Resolved at source (E9).} The printed Eq.~(87) ($e^{-2\alpha}$
suppression) yields masses uniformly high by $1.0005504$; the author's code
carries $e^{-(2+1/4\pi)\alpha}$, and the extra $e^{-\alpha/4\pi}$ is exactly
the missing class. Run verbatim, the code gives W at $0.17\sigma$, Z at
$1.25\sigma$, H at $0.05\sigma$. Repair: print the full exponent and derive
the $1/(4\pi)$ term. Two third-digit slips: the printed $g\approx0.6205$,
$g'\approx0.3469$ should be $0.6198$, $0.3471$ per the printed formulas (and
the printed pair is mutually inconsistent in $\sin^2\theta_W$).

\subsection{CKM and PMNS (Theorems 42--44)}
\label{sec:ckm}

From our independently generated quark masses:
$|V_{us}|=\sqrt{m_d/m_s}=0.223270$ (printed $0.2233$; $-1.54\sigma$) and
$|V_{cb}|=\tfrac23\bigl|\sqrt{m_s/m_b}-\sqrt{m_c/m_t}\bigr|=0.042632$
(printed $0.0426$; $+0.66\sigma$). Both survive substituting PDG central
masses. We note for completeness: $|V_{us}|$ is the credited
Gatto--Sartori--Tonin relation whose $-1.5\sigma$ residual is its known
leading-order ceiling; $|V_{ub}|$ is bounded ($0\le F\le1$) but not
numerically predicted; and PMNS is a qualitative large-angle argument whose
only numeric, $\Delta m^2_{31}/\Delta m^2_{21}=32.84$ (printed $32.8$),
follows from the neutrino spectrum.

\subsection{Cosmological constant (Theorems 57--58)}
\label{sec:lambda}

$\Lambda=3\exp\bigl(-(2+\zeta(3)/24)/\alpha\bigr)$ in Planck units (the
identity $\sigma_3\zeta(2)=\zeta(3)/24$ checks algebraically). Evaluated:
exponent $280.936$, $\Lambda=2.940\times10^{-122}$---the printed value under
\emph{both} $\alpha$ conventions; the E4 $\alpha$ gap moves $\Lambda$ by only
$0.017\%$, so this prediction does not inherit the E6 cancellation issue. We
record the structural fragility: with $e^{-281}$, a $0.4\%$ exponent error
moves $\Lambda$ by a factor of~3. $H_0=68.5$ is printed via the framework's
Friedmann sector and was not independently re-derived.

\subsection{Novel predictions (Part IV)}
\label{sec:partiv}

All arithmetic verifies: $|\Omega|=\alpha/2\pi=1.1614\times10^{-3}$,
$|g|=\alpha^2/4\pi^2=1.3489\times10^{-6}$, prefactors $\alpha^2/4\pi=
4.2376\times10^{-6}$ and $\alpha^3/8\pi^2=4.9216\times10^{-9}$, and all
fifteen entries of the five-configuration interferometer table (lab bench,
piezo-enhanced, ZARM free-fall tower, AION-10, AION-100; $\Delta\phi$,
$\delta\theta$, $\theta_{\mathrm{pol}}$). The phase wobble $\Delta\phi=
4.2\times10^{-6}$~rad for a tabletop Mach--Zehnder---four orders of magnitude
above the $10^{-10}$~rad/$\sqrt{\mathrm{Hz}}$ floor, with GR predicting
exactly zero---is the cleanest near-term falsifier of the framework. Two
interpretive flags: the identification of the measured LaAlO$_3$/SrTiO$_3$
Bloch-band QGT with the spacetime fiber QGT is asserted rather than derived;
and the dark-sector structural predictions ($w=-1$, flat rotation curves,
discrete $v_0$ clustering, no dark-matter particle) are qualitative in v5,
with no numerical $\lambda_\Omega$ yet printed for
$v_0^2=4\pi G\lambda_\Omega^2 n$.

\section{Machine-checked kernels (author's Lean 4 file)}
\label{sec:lean}

\texttt{tuft\_verify.lean} (self-contained, no Mathlib) formalizes six
structural kernels: charge quantization $\Rightarrow$ nontrivial holonomy;
domain indecomposability; uniqueness of classifying objects up to homotopy
equivalence (U(1) forcing onto Milnor's $\mathbb{CP}^\infty$ model);
completeness-as-theorem; and second-difference invariance of the affine
absorption convention. Checked under Lean~4.34.0: \textbf{all six compile and
check in $\sim$3~s, with no errors and no \texttt{sorry}}. The
\texttt{\#print axioms} audit is clean: four kernels depend on no axioms at
all; two touch only \texttt{propext}/\texttt{Quot.sound}. Milnor's theorem
correctly enters as a named hypothesis. One trivial erratum: as committed,
the file does not compile because the block-comment opener \texttt{/-} was
clobbered by a paste artifact; restoring it (one line) fixes the build. The
kernels verify the logical skeleton of the forcing chain; the physical
identifications remain outside the Lean scope, as the file itself is careful
to state.

\section{Master results table}

\begin{longtable}{@{}llll@{}}
\toprule
Claim & Source & Result & Status \\
\midrule
\endhead
UV completion & MQGT-SCF & reproduces; NGFP outside domain & E1 wording \\
Neutrino portal & MQGT-SCF & inherits TUFT Eq.~(136) & E2 resolved \\
T-3 constants & audit/TUFT & both identified on $\Sp7$/$\Sp9$, re-derived & E3 closed \\
$\alpha$ identity & Thm.~48 & exact; Wyler-identical; $6\times10^{-7}$ gap & E4 open \\
Pontryagin classes & T-1 script & corrected $5h^2$, $10h^4$ & E5 fixed \\
$g-2$ system & Eqs.~(144--148) & all digits reproduce; $\alpha$ caveat & E6 noted \\
Lepton spectrum & Thm.~34 & $n$-dependence exact; one const.\ gap & E7 repair stated \\
Quark spectrum & Thm.~37 & six masses, $\le10^{-6}$, $\le0.35\sigma$ & E8 clean pass \\
Boson spectrum & Thm.~35 & structure exact; $e^{-\alpha/4\pi}$ found & E9 resolved \\
CKM & Thm.~42 & $V_{us}$, $V_{cb}$ verify & E10 verified \\
Cosmological $\Lambda$ & Thm.~57/58 & $2.940\times10^{-122}$, both $\alpha$ & E11 verified \\
Interferometer table & Part IV & 15/15 entries & E12 verified \\
Lean kernels & author repo & 6/6 check, axiom-clean & E13 verified \\
\bottomrule
\end{longtable}

\section{Open items}

Two items remain genuinely open, both now stated narrowly enough to attack:
(i) the \emph{equipartition postulate} behind the framing numbers $\ell_7=56$
and $\ell_9=16$ in the universal phase---asserted, not derived; and (ii) the
$6.08\times10^{-7}$ gap between the derived and measured $\alpha$, which the
electron $g-2$ agreement partially rests upon (E6), and for which the
program's own scan shows gate-passing corrections are densely populated by
accidentals (E4). The discriminating experimental test is $a_\tau$.

\section{Reproducibility}

Every number in this report is rerunnable from the public repository
\texttt{github.com/Cbaird26/mqgt-scf-independent-verification}: the errata
document (\texttt{MQGT\_SCF\_ERRATA\_2026-09-17.md}, v4), the claims
inventory, the MQGT-SCF verification scripts
(\texttt{mqgt\_c6\_independent\_uv.py}, \texttt{mqgt\_t3\_independent.py},
\texttt{mqgt\_t3\_reverse\_search.py}, \texttt{mqgt\_t1\_independent.py},
\texttt{mqgt\_t1\_topological\_scan.py}), the TUFT verification script
(\texttt{tuft\_neutrino\_alpha\_verify.py}, sections [A]--[J]), and the
repaired Lean file (\texttt{tuft\_verify\_fixed.lean}).

\begin{thebibliography}{9}
\bibitem{nielsen2026} J.~L.~Nielsen, \textit{The Complex Hopf Fibration as
the Canonical Space for Gauge--Gravity Unification: The Field, Universal
Action, and Particle Spectrum}, Preprints.org 202604.0315.v5 (2026),
doi:10.20944/preprints202604.0315.v5.
\bibitem{nielsenrepo} J.~L.~Nielsen, \textit{Center for Topological Physics}
(code repository), \texttt{github.com/startigerjln/CenterforTopologicalPhysics},
retrieved 2026-09-17.
\bibitem{mqgt} C.~M.~Baird, MQGT-SCF program, canonical head
\texttt{mqgt-scf-science-public} (2026-09-14); independent verification
repository \texttt{github.com/Cbaird26/mqgt-scf-independent-verification}.
\bibitem{pdg} Particle Data Group, \textit{Review of Particle Physics}
(values as quoted in~\cite{nielsen2026}).
\bibitem{wp25} 2025 lattice-QCD White Paper muon $g-2$ determination, as
quoted in~\cite{nielsen2026}.
\bibitem{wyler} A.~Wyler, \textit{C.~R.~Acad.\ Sci.\ Paris} \textbf{A269}
(1969) 743; \textbf{A271} (1971) 186.
\end{thebibliography}

\end{document}
